\section*{Introduction}
Clinical AI systems can gain predictive information by combining many sensitive observations. That creates a practical question: how much of their useful performance survives when patients limit which data the system may process? We examine this question through sepsis prediction, comparing full-input reference models with the same models under simulated patient-specific consent restrictions. The comparison concerns permitted information use at inference time. It does not treat absence of consent as permission to collect data, and it does not evaluate removal of a patient's influence from an already trained model.

TRACER motivates the broader problem by combining continuous bedside video with clinical records for action recognition and longitudinal reasoning \citep{tracer}. Its authors report institutional review approval and restricted handling of protected information, while acknowledging consent and staff-surveillance questions for deployment. Their recordings are not publicly available. Our sepsis experiment therefore addresses a related permission-versus-performance question on public clinical data; it neither reproduces TRACER's tasks nor establishes comparable results on its private benchmark.

We frame this work as a sepsis benchmark study: a reusable protocol for evaluating specified permission interventions and the operating consequences of those interventions. The benchmark contributes paired comparisons and documented failure cases rather than a new learning algorithm or a claim to measure ethics as a single score.

The central empirical comparison holds the trained predictor, test patients, outcome definitions and alert thresholds fixed. It changes the permitted inputs through partial opt-out, denial of individual feature groups, complete physiological-input denial, and withdrawal during monitoring. Paired patient-level resampling estimates changes in discrimination, missed incidents, review demand and benchmark utility. Holding the decision rule fixed separates this comparison from our earlier sensitivity analysis, which selected a different empirical threshold for each consent scenario.

The mathematical analysis defines paired performance changes and a bounded-loss preservation criterion with a complete proof. It separates a statistical claim of non-inferiority at a declared tolerance from an ethical judgment about which information patients should permit. Supplementary analyses examine alert timing, finite review capacity and modeled oversight; they provide context for interpreting downstream behavior without changing the benchmark's central full-input-versus-constrained comparison.

\section*{Related work}
The PhysioNet challenge was designed to make early sepsis prediction comparable across cohorts, clinical variables and evaluation procedures \citep{reyna2020}. Its utility rewards prediction at different times relative to onset and penalizes false alarms. That makes it more informative than a single threshold-free ranking metric, but it still evaluates a prescribed prediction representation. Our once-per-patient review count and strict first-review window are additional estimands, not replacements for the challenge's original definition.

Clinical time-series benchmarks have established the value of comparing linear and recurrent models within a common preprocessing and evaluation framework \citep{harutyunyan2019}. Missingness is also informative: the fact that a measurement was ordered can predict outcomes, independently of the recorded value. GRU-D explicitly represents missingness and time intervals \citep{che2018}. Our GRU is a simpler baseline using causal summaries and observation counts; it does not implement learned decay or reproduce GRU-D's architecture. This distinction matters when interpreting the modest model search as evidence about a policy family rather than a comprehensive search for the best sepsis predictor.

Discrimination, probability quality and clinical usefulness also differ. Precision--recall analysis is particularly informative when positives are rare \citep{saito2015}; calibration studies emphasize that apparently strong ranking does not ensure reliable risk probabilities \citep{vancalster2019}. We therefore tune using validation average precision and report AUROC, average precision, Brier score and challenge utility. Class-weighted training can improve the ordering of cases while changing the scale of the predicted probabilities, so we do not interpret the raw weighted scores as calibrated clinical risks.

The external validation of the Epic Sepsis Model examined both predictive performance and the burden of alternative alerting strategies \citep{wong2021}. The TREWS adoption study examined whether clinicians evaluated alerts and how their interactions related to treatment timing \citep{trews}. These studies motivate our separation between an alert being emitted and a correct action being completed. They do not supply the active service-time distribution for our simulated reviewers; response delay includes workflow and availability as well as time spent on the assessment.

Dynamic consent treats permission as an ongoing interaction in which participants can manage preferences as data use evolves \citep{kaye2014}. We model one technical consequence of changing permissions: which inputs and derived state are available to a predictor. This supports controlled sensitivity experiments but omits comprehension, communication, preference formation and the institutional process through which consent is obtained.

Meaningful human control has also been formulated in terms of systems responding to relevant human reasons and outcomes remaining attributable to responsible people \citep{santoni2018}. A review node in an information graph captures only a routing requirement. It says neither that a reviewer has enough information nor that the reviewer has time, competence or authority to act. Our graph checks and stochastic reviewer model address distinct parts of that larger process, while prospective human evaluation remains outside the present study.
